% GENERATED FILE. Edit canonical lessons, then run npm run generate:books.
% canonical-source-hash: fnv1a64:bde0a413f0c301ff
% canonical-lessons: FR-C220-ponctuel-ponctuelle, FR-C220-etranger-etrangere, FR-C220-aimable, FR-C220-honnete, FR-C220-patient-patiente

\chapter{Punctual, Foreign, Friendly, Honest, Patient}
\label{ch:fr-punctual-foreign-friendly-honest-patient}
\hlchaptermodality{220}

\begin{chapteropening}
\textbf{By the end of this chapter:} \emph{I can say punctual, foreign, friendly, honest and patient.}

Complete the last lesson of chapter 220: I can say punctual, foreign, friendly, honest and patient.
\end{chapteropening}

\section[ponctuel / ponctuelle]{ponctuel / ponctuelle — punctual}

\label{lesson:FR-C220-ponctuel-ponctuelle}



\begin{quote}
Before the new one: say the French for different, then the French for the same.
\end{quote}

\subsection*{You'll want to know: ponctuel / ponctuelle}
\textbf{ponctuel / ponctuelle} — \textquotedblleft{}punctual\textquotedblright{}.

\begin{cousinweb}
\textbf{ponctuel} $\leftarrow$ Medieval Latin \emph{punctualis}, from \emph{punctus}, \textquotedblleft{}a point\textquotedblright{} — a learned borrowing. (EtymDB 2.1)
\end{cousinweb}

\begin{grammarlens}[title={What you've built}]
One more word for this chapter.
\end{grammarlens}

\subsection*{Your turn}
\begin{itemize}
\raggedright
  \item \emph{Say it:} \emph{ponctuel / ponctuelle}
  \item \emph{Say it:} \emph{ponctuel / ponctuelle}, once more
  \item \emph{Say it:} say \emph{ponctuel / ponctuelle}
  \item \emph{Recall:} say the French for different, then the French for the same, then say \emph{ponctuel / ponctuelle} again
\end{itemize}

\begin{tcolorbox}[breakable,colback=teal!4,colframe=teal!35!black,title={Before you move on}]
What does ponctuel / ponctuelle mean? (\textquotedblleft{}Punctual\textquotedblright{}.) Say it once more.
\end{tcolorbox}
\section[étranger / étrangère]{étranger / étrangère — foreign}

\label{lesson:FR-C220-etranger-etrangere}



\begin{quote}
Before the new one: say the French for the same, then the French for punctual.
\end{quote}

\subsection*{You'll want to know: étranger / étrangère}
\textbf{étranger / étrangère} — \textquotedblleft{}foreign\textquotedblright{}.

\begin{cousinweb}
\textbf{étranger} — built from \emph{étrange}, then meaning \textquotedblleft{}foreign\textquotedblright{}, with \emph{-ier}. (EtymDB 2.1)
\end{cousinweb}

\begin{grammarlens}[title={What you've built}]
One more word for this chapter.
\end{grammarlens}

\subsection*{Your turn}
\begin{itemize}
\raggedright
  \item \emph{Say it:} \emph{étranger / étrangère}
  \item \emph{Say it:} \emph{étranger / étrangère}, once more
  \item \emph{Say it:} say \emph{étranger / étrangère}
  \item \emph{Recall:} say the French for the same, then the French for punctual, then say \emph{étranger / étrangère} again
\end{itemize}

\begin{tcolorbox}[breakable,colback=teal!4,colframe=teal!35!black,title={Before you move on}]
What does étranger / étrangère mean? (\textquotedblleft{}Foreign\textquotedblright{}.) Say it once more.
\end{tcolorbox}
\section[aimable]{aimable — friendly}

\label{lesson:FR-C220-aimable}



\begin{quote}
Before the new one: say the French for punctual, then the French for foreign.
\end{quote}

\subsection*{You'll want to know: aimable}
\textbf{aimable} — \textquotedblleft{}friendly\textquotedblright{}.

\begin{cousinweb}
\textbf{aimable} $\leftarrow$ Latin \emph{amabilis}, \textquotedblleft{}lovable\textquotedblright{} (Lewis \& Short) — inherited, reshaped on \emph{aimer}.
\end{cousinweb}

\begin{grammarlens}[title={What you've built}]
One more word for this chapter.
\end{grammarlens}

\subsection*{Your turn}
\begin{itemize}
\raggedright
  \item \emph{Say it:} \emph{aimable}
  \item \emph{Say it:} \emph{aimable}, once more
  \item \emph{Say it:} say \emph{aimable}
  \item \emph{Recall:} say the French for punctual, then the French for foreign, then say \emph{aimable} again
\end{itemize}

\begin{tcolorbox}[breakable,colback=teal!4,colframe=teal!35!black,title={Before you move on}]
What does aimable mean? (\textquotedblleft{}Friendly\textquotedblright{}.) Say it once more.
\end{tcolorbox}
\section[honnête]{honnête — honest}

\label{lesson:FR-C220-honnete}



\begin{quote}
Before the new one: say the French for foreign, then the French for friendly.
\end{quote}

\subsection*{You'll want to know: honnête}
\textbf{honnête} — \textquotedblleft{}honest\textquotedblright{}.

\begin{cousinweb}
\textbf{honnête} $\leftarrow$ Old French \emph{honeste} $\leftarrow$ Latin \emph{honestus}, \textquotedblleft{}honourable\textquotedblright{}, from \emph{honor} (EtymDB 2.1) — inherited; the circumflex marks the lost \emph{s}.
\end{cousinweb}

\begin{grammarlens}[title={What you've built}]
One more word for this chapter.
\end{grammarlens}

\subsection*{Your turn}
\begin{itemize}
\raggedright
  \item \emph{Say it:} \emph{honnête}
  \item \emph{Say it:} \emph{honnête}, once more
  \item \emph{Say it:} say \emph{honnête}
  \item \emph{Recall:} say the French for foreign, then the French for friendly, then say \emph{honnête} again
\end{itemize}

\begin{tcolorbox}[breakable,colback=teal!4,colframe=teal!35!black,title={Before you move on}]
What does honnête mean? (\textquotedblleft{}Honest\textquotedblright{}.) Say it once more.
\end{tcolorbox}
\section[patient / patiente]{patient / patiente — patient}

\label{lesson:FR-C220-patient-patiente}



\begin{quote}
Before the new one: say the French for friendly, then the French for honest.
\end{quote}

\subsection*{You'll want to know: patient / patiente}
\textbf{patient / patiente} — \textquotedblleft{}patient\textquotedblright{}.

\begin{cousinweb}
\textbf{patient} $\leftarrow$ Latin \emph{patiens}, \emph{patientis}, \textquotedblleft{}enduring\textquotedblright{}, from \emph{pati}, \textquotedblleft{}to suffer\textquotedblright{} — a learned borrowing. (EtymDB 2.1)
\end{cousinweb}

\begin{grammarlens}[title={What you've built}]
One more word for this chapter.
\end{grammarlens}

\subsection*{Your turn}
\begin{itemize}
\raggedright
  \item \emph{Say it:} \emph{patient / patiente}
  \item \emph{Say it:} \emph{patient / patiente}, once more
  \item \emph{Say it:} say \emph{patient / patiente}
  \item \emph{Recall:} say the French for friendly, then the French for honest, then say \emph{patient / patiente} again
\end{itemize}

\begin{tcolorbox}[breakable,colback=teal!4,colframe=teal!35!black,title={Before you move on}]
What does patient / patiente mean? (\textquotedblleft{}Patient\textquotedblright{}.) Say it once more.
\end{tcolorbox}
